% =====================================================================
% sections/introduction.tex  --  input by report.tex.
%
% Word budget: ~700 words (counted toward the 3000-word limit).
%
% Written last-but-one, with the finished discussion in hand, so every setup
% here has a specific payoff there. Five beats, each setting up the next:
%   1. Working definition of adaptivity (Ashby / Di Paolo) -- the discussion
%      refers back to this by name, so it MUST be stated explicitly here.
%   2. HP background: biological homeostasis (Turrigiano) + the computational
%      substrate-level result (Williams & Noble) -> HP as cybernetic feedback.
%   3. The substrate-vs-evolvability tension: HP helps the substrate, yet
%      Williams found online HP did not improve evolved fitness.
%   4. Baldwin / genetic-assimilation frame: lifetime adaptation can guide
%      evolution and be assimilated; the conditions vary the lifetime adaptation;
%      the frozen-HP test probes assimilation directly.
%   5. Contribution statement: replication of Williams Ch.7 across four HP
%      conditions + four analyses, the central one being the frozen-HP test.
% =====================================================================

\section{Introduction}
\label{sec:intro}

An adaptive system, in the cybernetic tradition this project works within, is
one that maintains its essential variables within viable bounds despite
perturbation, by adjusting its own organisation across more than one timescale
\parencite{ashby1960}. Ashby's homeostat reorganises its internal connections
to return to equilibrium when pushed outside its stable region; Di Paolo's
evolved agents recover coordinated behaviour after their sensors are inverted,
adapting on a timescale slower than the behaviour itself
\parencite{dipaolo2000}. What both share is regulation: a slow process acting on
the parameters of a fast one to keep the fast one viable. This is the sense of
adaptivity used throughout, and the question the project asks is what happens
when such a regulatory process is placed in competition, and in collaboration,
with evolution.

The regulatory process in question is homeostatic plasticity (HP). In
biological neurons, homeostatic mechanisms hold mean firing rates near a
set-point by scaling synaptic strengths and adjusting excitability, stabilising
activity that Hebbian learning would otherwise drive to saturation or silence
\parencite{turrigiano1999}. The same idea transfers to artificial continuous-time
recurrent neural networks (CTRNNs): a plasticity rule that drives each neuron's
firing rate toward a viable range improves the network's behaviour as a
substrate, keeping neurons in the sensitive part of their activation function
rather than saturated at its rails \parencite{williamsnoble2007}. HP is, in
short, cybernetic negative feedback applied to a neural controller --- a slow
loop regulating a fast one, in exactly the form the definition above describes.

If HP improves the neural substrate, it is natural to expect it to improve the
\emph{evolvability} of that substrate: networks that stay in their sensitive
range should be easier for evolution to shape into useful controllers. Williams
tested this in his ball-catching experiments and found the relationship to be
more subtle than expected \parencite{williams2006}. HP applied as a
\emph{developmental} phase before each fitness evaluation improved evolved
performance, but HP left running \emph{during} behaviour did not, and could
make matters worse. The substrate-level benefit did not translate
straightforwardly into an evolutionary one. Why a process that demonstrably
helps the substrate should fail to help, or even hinder, when active during
evaluation is the empirical puzzle this project sets out to replicate and
probe.

The puzzle has a natural theoretical home in the Baldwin effect: the proposal
that adaptation acquired within an individual's lifetime can guide evolution
without being inherited, and that traits first acquired by lifetime adaptation
may, over generations, become encoded innately in the genotype --- genetically
assimilated \parencite{baldwin1896, hintonnowlan1987, mayley1996}. HP is a form
of lifetime adaptation, and the experimental conditions vary it directly:
development-only HP is precisely the Hinton-and-Nowlan setup, lifetime
adaptation before evaluation; HP during behaviour is the harder case, where the
adaptation and the fitness measurement are interleaved. The Baldwin frame turns
Williams's puzzle into a sharp question: has the behaviour evolved under HP been
genetically assimilated --- absorbed into the genotype so that it persists when
HP is removed --- or does it remain dependent on the plasticity that produced
it? Stolting et al.\ suggest a specific mechanism by which assimilation might
fail, namely that HP-during-behaviour can evolve dynamics that exist only in the
joint space of network state and plastic parameters and so cannot survive the
plasticity's removal \parencite{stolting2023}.

This project replicates Williams's Chapter~7 ball-catching evolvability
experiment across four HP conditions --- no HP, development-only, behaviour-only,
and both --- and extends it with four analyses he did not perform: a behavioural
trajectory inspection of evolved individuals, a per-neuron viable-range
diagnostic tracked across evolution, a population-level analysis of the search
itself, and, centrally, a frozen-HP test that removes plasticity from evolved
controllers to measure how much of their behaviour survives. The frozen-HP test
is the project's direct empirical answer to the assimilation question, and the
result --- that assimilation has not occurred in any HP condition --- organises
the discussion that follows.
